% !TEX root = ../main.tex
\chapter{Conclusions and Future Work}
\label{ch:con}

%--------------------------------------------------------------
\section{Conclusions}
%--------------------------------------------------------------

This project evaluated the effect of symlog-transformed losses on the stability and performance of imagination-based reinforcement learning agents. Inspired by DreamerV3~\citep{hafner2023dreamerv3}, we isolated and evaluated the contribution of symlog in a simplified setting.

We implemented a simplified Dreamer-like agent operating in raw state space without the RSSM, two-hot encoding, or KL balancing. We compared three configurations: model-free actor-critic (Condition~A), imagination-based actor-critic with standard MSE losses (Condition~B), and imagination-based actor-critic with symlog-transformed losses (Condition~C). These were evaluated on CartPole-v1 and LunarLander-v3.

The key findings are:

\begin{itemize}
    \item \textbf{Imagination-based training improves sample efficiency in benign settings.} In CartPole-v1, Condition~B reached near-optimal performance in fewer steps than Condition~A, supporting the premise that learning from simulated experience can improve sample efficiency.

    \item \textbf{Standard MSE losses can cause value function divergence in imagination-based agents.} In LunarLander-v3, where rewards span from $-100$ to $+100$, Condition~B failed to learn. The critic loss escalated to magnitudes exceeding $10^{15}$, leading to policy collapse with mean returns of approximately $-1000$.

    \item \textbf{Symlog-transformed losses prevent value divergence and reward prediction instability.} Condition~C maintained stable critic losses in both environments. The logarithmic compression of the symlog function bounds gradient magnitudes, preventing the feedback loop between prediction errors and large parameter updates.

    \item \textbf{Symlog provides stability independent of other architectural components.} The stabilizing effect of symlog was observed without other DreamerV3 components, indicating that symlog directly addresses value divergence in imagination-based learning.
\end{itemize}

This study provides an evaluation of the symlog transformation in isolation. By controlling for other components of DreamerV3, we show that symlog alone is sufficient to prevent value function divergence. This suggests that symlog-transformed losses are a useful design choice for agents training on simulated trajectories.

%--------------------------------------------------------------
\section{Future Work}
%--------------------------------------------------------------

The limitations identified in Chapter~\ref{ch:evaluation} suggest several directions for future research.

\subsection{Short-Term Extensions}

\textbf{Continuous action spaces.} A direct extension is to test the configurations in continuous action spaces, such as MountainCarContinuous-v0 or HalfCheetah-v4. Continuous control tasks feature different reward distributions and require policy gradient methods that may interact differently with the symlog transformation.

\textbf{Two-hot encoding.} DreamerV3 combines symlog with two-hot categorical encoding for value and reward predictions~\citep{hafner2023dreamerv3}. Future work could introduce two-hot encoding as a separate experimental configuration to analyze the individual and combined effects of symlog and two-hot encoding.

\textbf{Reward scaling experiments.} To characterize the sensitivity of imagination-based training to reward scale, experiments could vary reward scales by multiplicative factors (such as $100\times$ or $0.01\times$). This would show the thresholds at which value divergence occurs under standard MSE losses.

\subsection{Long-Term Investigations}

\textbf{Recurrent latent dynamics.} Replacing the MLP-based world model with a Recurrent State-Space Model (RSSM) would bring the architecture closer to DreamerV3. This would allow investigation of whether symlog remains necessary when using latent-space world models.

\textbf{Complex environments.} Evaluating the agents on pixel-based observations or tasks with longer time horizons, such as Atari games, would test the scalability of the approach.

\textbf{Alternative normalization techniques.} Future work could compare symlog with other normalization methods, such as PopArt (adaptive normalization of targets) or running-mean normalization, to establish performance trade-offs.

%--------------------------------------------------------------
\section{Summary}
%--------------------------------------------------------------

This chapter summarized the findings, highlighting the role of symlog-transformed losses in preventing value divergence during imagination-based reinforcement learning. The results show that symlog is an effective stabilization mechanism in isolation. Future work includes evaluating continuous action spaces, alternative encoding methods, and more complex environments.